\frontsection{Общие требования к отчёту и рекомендации преподавателю}
\subsection*{Структура отчёта}
Для каждой работы укажите ФИО, ИТМО ID, номер и название. Далее приведите цель, описание исходных данных, протокол разбиения, исходную модель, план опытов, результаты и выводы. Полные таблицы данных и многотысячные листинги не требуется вставлять в PDF: они передаются отдельными файлами. В тексте должны оставаться сведения, необходимые для понимания и проверки вывода.

\begin{enumerate}
\item \textbf{Постановка.} Что предсказывается, какие данные доступны и какая ошибка существенна.
\item \textbf{Протокол.} Как получены части данных, что обучалось только на training, по какой метрике выбирались настройки.
\item \textbf{Модель.} Принцип работы, архитектура или алгоритм, параметры опыта и мотивация выбора.
\item \textbf{Результаты.} Таблица всех запланированных опытов, необходимые графики, разбор ошибок.
\item \textbf{Вывод.} Наблюдение, возможное объяснение, ограничение и следующий проверяемый шаг.
\end{enumerate}

\subsection*{Критерии оценивания}
\begin{table}[H]\centering\small
\begin{tabularx}{\textwidth}{@{}p{0.25\textwidth}Y@{}}
\toprule
Критерий & Наблюдаемый признак выполнения \\
\midrule
Данные & Известны признаки и цель; нет утечки, перепутанных выборок и неверных меток. \\
Модель & Студент объясняет алгоритм, размерности, потери и изменяемые настройки. \\
Эксперимент & Сравнения воспроизводимы; сохранены разбиения, seed, версии и результаты. \\
Анализ & Показаны ошибки и ограничения; выводы следуют из полученных результатов. \\
Самостоятельность & Есть обоснованный выбор, исследование параметров и собственная архитектура в работе №\,3. \\
\bottomrule
\end{tabularx}\end{table}
Баллы, часы, минимальные пороги и правила пересдачи определяются рабочей программой дисциплины. Не следует автоматически оценивать работу только по максимальной метрике: утечка может дать высокий результат при неправильном эксперименте. Напротив, корректно исследованный неудачный вариант может содержать содержательный учебный результат.

\subsection*{Как организовать занятие}
До начала курса подготовьте проверенные версии файлов и описаний, уточните названия служебных столбцов, распределите доступные вычислительные ресурсы. Для работы №\,3 отдельно оцените память исходной архитектуры. После пробного выполнения студентом с заявленной подготовкой уточните трудоёмкость; произвольное количество часов в пособии не назначено.

Полезно проводить короткий допуск: студент объясняет роли файлов, рисует разбиение и вычисляет один пример метрики или размерности. В работе №\,1 проверяется логика сравнения, в №\,2 --- согласованная подготовка признаков и шкала цены, в №\,3 --- режимы сети и самостоятельность архитектурного решения.

Для первого курса машинного обучения полезно разделить занятие на три прохода: сначала разобрать данные и получить один прогноз, затем объяснить его и метрику, после этого переходить к таблице опытов. На первом проходе не требуется одновременно менять все настройки. Обзор современных табличных моделей и опыт с ансамблем являются дополнительными и не входят в обязательный минимум.

Стартовые ноутбуки содержат рабочие примеры. На защите предложите студенту объяснить конкретную строку, предсказать эффект изменения параметра и сопоставить это ожидание со своим результатом. Не требуйте скрытого эталонного значения accuracy или MAE, которое зависит от версии данных и протокола.

\subsection*{Дополнительные исследования}
После обязательного минимума можно исследовать другой способ извлечения признаков активности, содержательную обработку отсутствующих характеристик домов, аугментации изображений или частичный fine-tuning. Каждое расширение должно иметь собственную гипотезу и фиксированные условия; увеличение числа опытов без вопроса не делает исследование более убедительным.
